\begin{abstract}
Classical Tucker decomposition is built on multilinear low-rank structure. However, when an observed tensor is generated by nonlinear mechanisms, saturated responses, complex higher-order interactions, or heteroscedastic perturbations, the purely multilinear assumption often leads to systematic mismatch. To address this issue, this paper proposes Nonlinear Tucker Decomposition with Polynomial Links (NTD-PL). The model retains a shared Tucker latent tensor as the low-rank backbone and imposes an explicit polynomial link on top of it, thereby incorporating standard Tucker into a unified framework as a linear special case and enabling a controlled extension toward higher-order interactions through only a small number of link coefficients.

On the theoretical side, we clarify the nesting relation between NTD-PL and Tucker, establish exact representation results under polynomial generation, and derive finite-order truncation approximation bounds under analytic generation. We also formulate a unified masked learning objective and, together with Tucker initialization, gradient updates, and coefficient updates for the polynomial link, obtain a practical alternating optimization method. Experimental results show that, on linear source data, NTD-PL almost coincides with Tucker and does not create artificial gains; under a unified nonlinear residual protocol, NTD-PL achieves stable improvements in both structure-matching and finite-order approximation settings; and on real hyperspectral reconstruction and random-missing completion tasks, NTD-PL consistently outperforms Tucker under the same parameter budget. Further mechanism studies show that these gains can be attributed neither to simply adding parameters nor to applying a single shared polynomial correction to the Tucker output.

Overall, this paper provides an explicit, low-dimensional, and analyzable structured approach to enhancing nonlinear generation mechanisms while preserving the Tucker backbone.
\end{abstract}
